\section{Theory and Analytical Design}

\subsection{Wavelength and first-order length}
Using the design value $c\approx\SI{3e8}{m/s}$,
\begin{equation}
\lambda_0=\frac{c}{f_0}
=\frac{\SI{3e8}{m/s}}{\SI{1e9}{Hz}}
=\SI{0.300}{m}=\SI{300}{mm}.
\label{eq:wavelength}
\end{equation}
The first-order half-wave length is therefore
\begin{equation}
L_0=\frac{\lambda_0}{2}=\SI{150}{mm}.
\label{eq:halfwave-length}
\end{equation}
This is an analytical starting point rather than a guarantee of resonance for a
finite-diameter, gapped, port-excited model.

\subsection{Canonical resonant-dipole pattern}
For a thin, center-fed resonant half-wave dipole aligned with the $z$-axis, the normalized
far-field magnitude pattern is approximated by~\cite{stutzman2012}
\begin{equation}
F(\theta)=\left|
\frac{\cos\!\left(\frac{\pi}{2}\cos\theta\right)}{\sin\theta}
\right|.
\label{eq:half-wave-pattern}
\end{equation}
The maximum occurs at $\theta=90^\circ$ and the axial nulls occur at
$\theta=0^\circ$ and $180^\circ$. Numerical integration of the power pattern gives
\begin{equation}
D=1.6409=\SI{2.15}{dBi},
\end{equation}
and the $-\SI{3}{dB}$ points give an HPBW of $78.08^\circ$.

\subsection{Reflection coefficient and return loss}
For an antenna impedance $Z_A$ referenced to the CST port impedance $Z_0$,
\begin{equation}
\Gamma=\frac{Z_A-Z_0}{Z_A+Z_0},
\qquad
S_{11,\mathrm{dB}}=20\log_{10}|\Gamma|.
\label{eq:s11}
\end{equation}
Return loss is conventionally positive:
\begin{equation}
\mathrm{RL}=-20\log_{10}|\Gamma|=-S_{11,\mathrm{dB}}.
\label{eq:return-loss}
\end{equation}
Thus the final plotted value $S_{11}=-\SI{47.01186}{dB}$ corresponds to a return
loss of \SI{47.01186}{dB}, specifically under the \SI{73}{\ohm} reference used in the
CST port.

\subsection{Independent analytical calculation}
The script in Appendix~\ref{app:analytical-verification} evaluates
Eq.~\eqref{eq:half-wave-pattern}, finds the half-power crossings, integrates beam solid
angle, and records the comparison values independently of CST.
